\subsubsection{Jones Commutator Scalar and Achievability of Delinking Cost: Single-Parameter Fingerprint of the Index Gap and R\'enyi Flatness}\label{subsubsec:op-algebra-jones-scalar-twirl-cost-renyi-flatness}
This section, building on the index gap framework of Proposition \ref{prop:op-algebra-index-submultiplicativity-commuting-square}--Corollary \ref{cor:op-algebra-cmi-lower-bound-by-index-gap}, further shows that:
the index gap not only provides a lower bound on conditional mutual information, but can also be equivalently read as a \emph{single scalar obstruction} from the Jones projection sandwich;
moreover, in the caliber where the relative commutant is finite-dimensional and the index ratio is an integer, this lower bound is achievable---there exists a relative commutant twirl driven by finite randomness that achieves delinking (conditional independence) in the information-theoretic sense, with the required randomness cost exactly equal to the index gap.
Finally, a parallel ``R\'enyi flatness'' conclusion is given: the worst-case information loss of finite-index conditional expectations is exactly capped and achievable by the same constant \(\log\mathrm{Ind}(E)\) across the entire R\'enyi family.

\begin{proposition}[Jones--commutator scalar of the index gap: single-parameter obstruction of the projection sandwich]\label{prop:op-algebra-jones-scalar-equals-exp-minus-gap}
Following the setup of Corollary \ref{cor:op-algebra-cmi-lower-bound-by-index-gap}:
let \(\mathcal{A}\) be a finite von Neumann algebra,
\(\mathcal{B}_1,\mathcal{B}_2\subseteq\mathcal{A}\) subalgebras,
\(\mathcal{B}_0:=\mathcal{B}_1\cap\mathcal{B}_2\), with finite-index conditional expectations
\(
E_i:\mathcal{A}\to\mathcal{B}_i
\) (\(i=1,2\)) and
\(
E_{12}:\mathcal{A}\to\mathcal{B}_0
\).
Write the index gap as
\[
\Gamma:=\log\frac{\mathrm{Ind}(E_1)\mathrm{Ind}(E_2)}{\mathrm{Ind}(E_{12})}\ \ge\ 0.
\]
In the Jones basic construction, let \(e_1,e_2,e_{12}\) be the Jones projections corresponding to \(E_1,E_2,E_{12}\).
If it is further assumed that the centers of the compressed algebras \(e_1\langle\mathcal{A},e_{12}\rangle e_1\) and \(e_2\langle\mathcal{A},e_{12}\rangle e_2\) are both one-dimensional (equivalently: the relative commutant yields only a scalar in this compressed caliber), then there exists a unique scalar \(\lambda\in(0,1]\) such that
\[
\boxed{\;
e_1e_2e_1=\lambda\,e_1,\qquad e_2e_1e_2=\lambda\,e_2\;}
\]
and this scalar satisfies
\[
\boxed{\;
\lambda=\frac{\mathrm{Ind}(E_{12})}{\mathrm{Ind}(E_1)\mathrm{Ind}(E_2)}
=\exp(-\Gamma).\;}
\]
Therefore, the index gap \(\Gamma\) can be equivalently stated as the commutator scalar obstruction of the Jones projection product; it is strictly additive under tensorization (by the multiplicativity of the index and logarithmic transformation).
\end{proposition}
\begin{proof}[Proof sketch]
The operator \(e_1e_2e_1\) belongs to the commutative part of the compressed algebra \(e_1\langle\mathcal{A},e_{12}\rangle e_1\); under the one-dimensional center assumption, it must be a scalar multiple of \(e_1\), so there exists \(\lambda\) with \(e_1e_2e_1=\lambda e_1\).
The same argument applies to \(e_2e_1e_2\).
\par
Taking the canonical trace on both sides and using the standard relation between traces and indices in the Jones basic construction (e.g., the analogue of \(\tau(e_i)=\mathrm{Ind}(E_i)^{-1}\)) yields
\(
\lambda=\mathrm{Ind}(E_{12})/(\mathrm{Ind}(E_1)\mathrm{Ind}(E_2))
\).
The last identity follows immediately from the definition of \(\Gamma\). \qedhere
\end{proof}

\begin{proposition}[Achievability of delinking cost: relative commutant twirl achieves the index gap with exact randomness]\label{prop:op-algebra-twirl-achieves-index-gap}
Following the setup of Proposition \ref{prop:op-algebra-jones-scalar-equals-exp-minus-gap}, further assume that the relative commutant \((\mathcal{B}_0'\cap\mathcal{A})\) is finite-dimensional.
If the index ratio
\[
\frac{\mathrm{Ind}(E_1)\mathrm{Ind}(E_2)}{\mathrm{Ind}(E_{12})}
\]
is an integer, then there exists a random unitary completely positive map (twirl) driven by finite randomness
\[
\mathcal{T}(x):=\sum_{g=1}^{G}p_g\,u_g^\ast x\,u_g,\qquad x\in\mathcal{A},
\]
where \(u_g\in(\mathcal{B}_0'\cap\mathcal{A})\) are finitely many unitaries, \((p_g)_{g=1}^{G}\) is a probability distribution, satisfying:
\begin{enumerate}
  \item \(\mathcal{T}\) preserves a given faithful normal state \(\varphi\) (which can be chosen to be invariant under conjugation by \(u_g\));
  \item Under \(\varphi\circ\mathcal{T}\),
  \[
  \boxed{\;
  I_{\varphi\circ\mathcal{T}}(\mathcal{B}_1:\mathcal{B}_2\mid\mathcal{B}_0)=0, \;}
  \]
  i.e., the two sources are strictly delinked conditionally on the cross-visible factor \(\mathcal{B}_0\);
  \item The minimum Shannon entropy of the required randomness exactly equals the index gap:
  \[
  \boxed{\;
  \inf H(R)=\Gamma, \;}
  \]
  where \(R\) is the external random variable implementing \(\mathcal{T}\) (the classical register selecting \(g\)).
\end{enumerate}
Therefore, the index gap lower bound given by Corollary \ref{cor:op-algebra-cmi-lower-bound-by-index-gap} is achievable in this finite-depth/finite-commutant caliber: the index gap is not only the lower bound on the irreducible linking residual, but also the exact value of the minimum randomness required for ``proper denoising.''
\end{proposition}
\begin{proof}[Proof sketch]
When \((\mathcal{B}_0'\cap\mathcal{A})\) is finite-dimensional and the index ratio is an integer, one can select a finite group (or finite system of units) action on the relative commutant and take \(\mathcal{T}\) as its conjugation average.
This twirl is bimodular on \(\mathcal{B}_0\) and can be constructed to be \(\varphi\)-preserving, thereby achieving complete decoherence of the ``cross-interface'' in the information-theoretic sense; the delinking conclusion corresponds to the characterization of \(I_\varphi=0\) in the minimal commuting square case (Corollary \ref{cor:op-algebra-cmi-lower-bound-by-index-gap}).
\par
For the randomness cost: on one hand, the structural lower bound \(I_\varphi(\mathcal{B}_1:\mathcal{B}_2\mid\mathcal{B}_0)\ge\Gamma\) from Corollary \ref{cor:op-algebra-cmi-lower-bound-by-index-gap}, combined with the chain rule that ``random unitary averaging can only reduce conditional mutual information at the expense of classical entropy,'' gives that any mechanism achieving delinking must satisfy \(H(R)\ge\Gamma\).
On the other hand, taking a uniform twirl of size \(G=\exp(\Gamma)\) (\(p_g\equiv 1/G\)) yields \(H(R)=\log G=\Gamma\) and achieves \(I=0\);
hence the lower bound is achieved and the minimum equals \(\Gamma\). \qedhere
\end{proof}

\begin{theorem}[Optimal ``delinking curve'' of the index gap: strong duality of the randomness budget]\label{thm:op-algebra-optimal-delinking-curve}
Following the setup of Proposition \ref{prop:op-algebra-jones-scalar-equals-exp-minus-gap}, write the index gap as
\[
\Gamma:=\log\frac{\mathrm{Ind}(E_1)\mathrm{Ind}(E_2)}{\mathrm{Ind}(E_{12})}\ \ge\ 0.
\]
Consider a family of delinking mechanisms \(\mathcal{T}\) (normal u.c.p.\ maps; understood as random unitary averages or their combinations) that are allowed to introduce an external random register \(R\), with protocol complexity constrained by the Shannon randomness budget \(H(R)\le h\).
Write the conditional mutual information leakage relative to \(\varphi\circ\mathcal{T}\) as
\[
\mathcal{L}(\mathcal{T})
:=
I_{\varphi\circ\mathcal{T}}(\mathcal{B}_1:\mathcal{B}_2\mid\mathcal{B}_0).
\]
Then an unbreachable lower bound exists:
\[
\boxed{\;
\inf_{\mathcal{T}:\ H(R)\le h}\ \mathcal{L}(\mathcal{T})
\ \ge\ \max\{\Gamma-h,\,0\}.\;}
\]
Moreover, in the finite-depth/finite-relative-commutant case where the index ratio can be quantized to an integer, this lower bound is achievable: there exists a family \(\mathcal{T}_h\) such that
\[
\boxed{\;\mathcal{L}(\mathcal{T}_h)=\max\{\Gamma-h,\,0\}.\;}
\]
Therefore, \(\Gamma\) is simultaneously the generator of the ``minimum randomness cost'' and the ``minimum leakage curve,'' providing the strict strong duality optimality of delinking.
\end{theorem}
\begin{proof}[Proof sketch]
Incorporating the external random register into the second interface, the chain rule gives
\[
I_{\varphi\circ\mathcal{T}}(\mathcal{B}_1:\mathcal{B}_2,R\mid\mathcal{B}_0)
=\mathcal{L}(\mathcal{T})+I_{\varphi\circ\mathcal{T}}(\mathcal{B}_1:R\mid\mathcal{B}_0,\mathcal{B}_2)
\le \mathcal{L}(\mathcal{T})+H(R).
\]
On the other hand, in the finite-index caliber of this paper, the index gap provides a structural lower bound on the conditional mutual information of the ``extended interface'' (an analogous application of Corollary \ref{cor:op-algebra-cmi-lower-bound-by-index-gap}):
\(
I_{\varphi\circ\mathcal{T}}(\mathcal{B}_1:\mathcal{B}_2,R\mid\mathcal{B}_0)\ge\Gamma
\).
Combining yields \(\mathcal{L}(\mathcal{T})\ge \Gamma-H(R)\), which together with \(\mathcal{L}(\mathcal{T})\ge 0\) gives the stated lower bound.
\par
Achievability comes from the existence of graded conjugation average families on the finite relative commutant: when the index ratio is an integer, one can select nested finite conjugacy classes/subgroups so that ``entropy already paid'' progressively cancels the gap, pushing the residual leakage to \(\Gamma-h\); when \(h\ge\Gamma\), this degenerates to the complete delinking of Proposition \ref{prop:op-algebra-twirl-achieves-index-gap}. \qedhere
\end{proof}

\begin{theorem}[Multi-round strong converse theorem: strong form of linear cumulative leakage and critical phase transition]\label{thm:op-algebra-multi-round-strong-converse-linear-leakage}
Following the setup of Theorem \ref{thm:op-algebra-optimal-delinking-curve}, consider an \(n\)-round tensorized system with any protocol family \(\mathcal{T}^{(n)}\) (allowed to use a total randomness register \(R^{(n)}\) satisfying \(H(R^{(n)})\le nh\)).
Then the strong converse lower bound holds:
\[
\boxed{\;
I\bigl(\mathcal{B}_1^{\otimes n}:\mathcal{B}_2^{\otimes n}\mid \mathcal{B}_0^{\otimes n}\bigr)_{\varphi^{\otimes n}\circ\mathcal{T}^{(n)}}
\ \ge\ n\max\{\Gamma-h,\,0\}.\;}
\]
In particular, if \(h<\Gamma\), then leakage grows linearly at rate at least \(\Gamma-h\); if \(h\ge \Gamma\), then in the achievable case there exist protocols that push leakage to \(0\) (exactly or to arbitrary approximation).
Hence delinking exhibits a single critical point \(h=\Gamma\), with a strict phase transition structure.
\end{theorem}
\begin{proof}[Proof sketch]
The index gap is strictly additive under tensorization (multiplicativity of the logarithmic index), so the \(n\)-round system has gap \(n\Gamma\), while the randomness budget satisfies \(H(R^{(n)})\le nh\).
Applying Theorem \ref{thm:op-algebra-optimal-delinking-curve} to the tensorized system gives
\(
I\ge \max\{n\Gamma-H(R^{(n)}),0\}\ge n\max\{\Gamma-h,0\}
\).
The achievability direction is provided by per-round independent realization \(\mathcal{T}_h^{\otimes n}\) (in the achievable caliber). \qedhere
\end{proof}

\begin{corollary}[Noncommutative R\'enyi flatness: worst-case R\'enyi loss of finite-index horizons is \texorpdfstring{$\alpha$}{alpha}-invariant]\label{cor:op-algebra-renyi-flatness-sup-equals-log-index}
Let \(\mathcal{B}\subseteq\mathcal{A}\) be a finite von Neumann algebra inclusion, and \(E:\mathcal{A}\to\mathcal{B}\) a faithful normal finite-index conditional expectation.
For any \(\alpha\in[1,\infty]\), let \(D_\alpha\) be the sandwiched R\'enyi relative entropy, and define the R\'enyi horizon loss
\[
L_\alpha(\rho):=D_\alpha(\rho\|\rho\!\circ\!E).
\]
Then the uniform extremal formula
\[
\boxed{\;
\sup_{\rho\in\mathcal{S}(\mathcal{A})}L_\alpha(\rho)=\log\mathrm{Ind}(E)\qquad(\forall \alpha\in[1,\infty])\;}
\]
holds.
Therefore, the ``worst-case information recovery'' is strictly flat across the Shannon, R\'enyi, and min-entropy full spectrum, exactly capped and achievable by the same index constant.
\end{corollary}
\begin{proof}[Proof sketch]
By the Pimsner--Popa inequality, \(\rho\le \mathrm{Ind}(E)\,(\rho\!\circ\!E)\), hence for all \(\alpha\in[1,\infty]\),
\(
L_\alpha(\rho)\le \log\mathrm{Ind}(E)
\)
(consistent with the cap of Corollary \ref{cor:op-algebra-renyi-loss-spectrum-capped-by-index}).
\par
For achievability: by Proposition \ref{prop:op-algebra-dmax-capacity-equals-log-index}, \(\sup_\rho L_\infty(\rho)=\log\mathrm{Ind}(E)\).
In the finite-dimensional commutative caliber, the single-point state on the maximum fiber simultaneously saturates all R\'enyi orders (Corollary \ref{cor:op-algebra-commutative-dmax-extremizers-maxfiber}), yielding achievability of the lower bound for any fixed \(\alpha\).
In the general case, near-extremal states can be constructed via the Jones tower/basic construction: for any \(\varepsilon>0\), there exists \(\rho_\varepsilon\) with \(L_\infty(\rho_\varepsilon)\ge \log\mathrm{Ind}(E)-\varepsilon\), and using the monotonicity and continuous approximation of \(\alpha\mapsto L_\alpha(\rho)\), the lower bound is pushed to \(\log\mathrm{Ind}(E)\).
Hence the supremum takes this value for each \(\alpha\). \qedhere
\end{proof}

\begin{corollary}[Externalized algebraic curvature: second-order response of the index gap controlled by Jones scalar curvature]\label{cor:op-algebra-index-gap-curvature-second-variation}
Under the setup of Proposition \ref{prop:op-algebra-jones-scalar-equals-exp-minus-gap}, let \(\lambda\in(0,1]\) be the Jones sandwich scalar and write \(\Gamma=-\log\lambda\).
Suppose there exists a differentiable one-parameter perturbation family (preserving \(\mathcal{B}_0\) and finite-index feasibility) such that \(\lambda=\lambda(\tau)\), \(\Gamma=\Gamma(\tau)\) are twice differentiable.
Then for any \(\tau\), the identity
\[
\boxed{\;
\Gamma''(\tau)=\bigl(-\log\lambda(\tau)\bigr)''=\frac{\lambda'(\tau)^2}{\lambda(\tau)^2}-\frac{\lambda''(\tau)}{\lambda(\tau)}\;}
\]
holds.
In particular, if \(\tau=0\) is a local minimum of the index gap (equivalently a local maximum of \(\lambda(\tau)\)), then \(\lambda'(0)=0\) and \(\lambda''(0)\le 0\), hence
\[
\boxed{\;\Gamma''(0)\ge 0.\;}
\]
Therefore, near a minimal commuting square, the index gap exhibits a strict ``curvature-type'' obstruction: if a first-order perturbation does not increase the gap, then the second order must be nonnegative.
\end{corollary}
\begin{proof}
The identity follows by directly computing the second derivative of \(\Gamma(\tau)=-\log\lambda(\tau)\).
The derivative sign conditions at a local minimum/maximum are elementary calculus. \qedhere
\end{proof}

\begin{proposition}[Minimal sufficiency gap of the horizon interface: unified lower bound for recoverability error and index gap]\label{prop:op-algebra-minimal-sufficiency-gap-by-index}
Let \(\mathcal{B}_0\subset \mathcal{B}_1,\mathcal{B}_2\subset \mathcal{A}\) be a subalgebra chain in a finite von Neumann algebra, with faithful normal finite-index conditional expectations
\(
E_i:\mathcal{A}\to\mathcal{B}_i
\) (\(i=0,1,2\)).
Write the index gap as
\[
\Gamma:=\log\frac{\mathrm{Ind}(E_1)\mathrm{Ind}(E_2)}{\mathrm{Ind}(E_0)}\ \ge\ 0.
\]
For any recovery mechanism (normal CPTP) \(R\) (taking \(\mathcal{B}_1\vee\mathcal{B}_2\) as the full visible input and outputting a state on \(\mathcal{A}\)), define the worst-case relative entropy recovery error
\[
\epsilon_R:=\sup_{\rho\in\mathcal{S}(\mathcal{A})}
D\Bigl(\rho\ \Big\|\ R\bigl(\rho\!\upharpoonright_{\mathcal{B}_1\vee\mathcal{B}_2}\bigr)\Bigr),
\]
where \(D\) is the Umegaki relative entropy and \(\rho\!\upharpoonright_{\mathcal{B}_1\vee\mathcal{B}_2}\) denotes the restriction of \(\rho\) to the visible composite algebra.
Then an ineliminable lower bound exists:
\[
\boxed{\;
\inf_R\ \epsilon_R\ \ge\ \Gamma.\;}
\]
Equality \(\epsilon_R=0\) (exact recovery in the Petz sufficiency sense, cf.\ \cite{Petz1986SufficientSubalgebrasRelativeEntropy,Petz2003MonotonicityRelativeEntropyRevisited}) holds if and only if \(\Gamma=0\) (index saturation / minimal commuting square caliber).
Therefore, the index gap \(\Gamma\) simultaneously serves as the lower bound on the randomness cost for delinking (Proposition \ref{prop:op-algebra-twirl-achieves-index-gap}) and as the unified irreversibility potential expressing ``the composite visible interface is insufficient to fully recover the ambient algebra state.''
\end{proposition}
\begin{proof}[Proof sketch]
If an exact recovery \(\epsilon_R=0\) exists, then \(\rho\) can be losslessly recovered from its restriction to \(\mathcal{B}_1\vee\mathcal{B}_2\) via \(R\), while maintaining relative entropy non-increase with equality for all \(\rho\); by Petz's sufficiency theorem, this forces the corresponding visible subalgebra to be sufficient for \(\mathcal{A}\) in the relative entropy sense, which in the finite-index context of this paper is equivalent to index saturation (\(\Gamma=0\)).
Conversely, when \(\Gamma=0\), the minimal commuting square caliber admits a chain-decomposed state-preserving conditional expectation and a Petz recovery implementation, allowing exact recovery with \(\epsilon_R=0\).
\par
For general \(\Gamma>0\), the index gap provides a structural irreversibility lower bound: no recovery depending only on \(\mathcal{B}_1\vee\mathcal{B}_2\) can simultaneously push the relative entropy error below \(\Gamma\) for all states; otherwise it would contradict the index gap lower bound of Corollary \ref{cor:op-algebra-cmi-lower-bound-by-index-gap} and its sufficiency characterization. \qedhere
\end{proof}

\endinput
